\documentclass[11pt,a4paper]{article}

\usepackage[utf8]{inputenc}
\usepackage[T1]{fontenc}
\usepackage[slovene,provide=*]{babel}
\usepackage[margin=2.5cm]{geometry}
\usepackage{enumitem}
\usepackage{amsmath}
\usepackage{titlesec}
\usepackage{xcolor}
\usepackage[colorlinks=true,linkcolor=black,urlcolor=blue]{hyperref}

\definecolor{scred}{RGB}{150,20,30}

\titleformat{\section}{\Large\bfseries\color{scred}}{\thesection}{1em}{}
\titleformat{\subsection}{\large\bfseries}{\thesubsection}{1em}{}

\setlist[itemize]{topsep=2pt,itemsep=2pt,parsep=0pt}

\title{\textbf{Matematične seminarske in raziskovalne naloge}\\[4pt]\large Navodila in predlogi tem za gimnazijo}
\author{Šolski center Krško--Sevnica}
\date{}

\begin{document}
\maketitle
\tableofcontents
\newpage

% =====================================================================
\section{Raziskovalna naloga}
% =====================================================================

Raziskovalna naloga je namenjena dijakom, ki si želijo večjega izziva: gre za obsežnejšo, samostojno nalogo (10--15 strani), v primerjavi s krajšo seminarsko nalogo (do 3 strani), ki je predstavljena v razdelku~\ref{sec:seminarska}.

\subsection{Navodila za izdelavo}

\subsubsection*{Struktura naloge}

\begin{enumerate}
  \item \textbf{Naslovna stran} -- ime šole, predmet (matematika), naslov naloge, ime in razred dijaka, mentor, kraj in datum.
  \item \textbf{Kazalo vsebine} -- avtomatsko generirano kazalo s stranmi.
  \item \textbf{Uvod} (pribl.\ pol strani do ene strani) -- zakaj je dijak izbral to temo, kaj je namen/cilj naloge (raziskovalno vprašanje), kratek pregled vsebine.
  \item \textbf{Osrednji (jedrni) del} -- glavnina naloge:
  \begin{itemize}
    \item teoretična osnova: definicije, izreki, dokazi (prilagojeni gimnazijski ravni);
    \item lastni rešeni primeri -- ne le prepisana teorija, ampak dijakov lasten prikaz razumevanja;
    \item po možnosti praktični/uporabni del: povezava s stroko, vsakdanjim življenjem ali drugim predmetom;
    \item grafi, skice, slike (GeoGebra, Desmos, Excel), vse oštevilčene in poimenovane;
    \item če je smiselno, kratek zgodovinski oris pojma ali matematika, ki ga je odkril.
  \end{itemize}
  \item \textbf{Zaključek} -- strnjen povzetek ugotovitev, lasten razmislek (kaj je dijak izvedel, kje so bile težave, kaj bi raziskal naprej).
  \item \textbf{Viri in literatura} -- enoten citatni slog (npr.\ APA), ločeno knjige in spletni viri (z datumom dostopa); vsak prevzet podatek, slika ali trditev ima naveden vir.
  \item \textbf{Priloge} (po potrebi) -- dodatni izračuni, izpeljave, podatkovne tabele, izvorna koda.
\end{enumerate}

\subsubsection*{Oblikovanje}

\begin{itemize}
  \item Obseg: običajno 10--15 strani vsebine (brez naslovnice, kazala in prilog).
  \item Pisava: Times New Roman 12~pt ali Arial 11~pt, razmik vrstic 1,5.
  \item Enačbe urejene z urejevalnikom enačb (Word Equation, \LaTeX) -- ne kot slike, če se le da.
  \item Strani oštevilčene; slike in tabele naslovljene ter oštevilčene (npr.\ ,,Slika 3: \dots{}``).
\end{itemize}

\subsubsection*{Nasveti za kakovostno nalogo}

\begin{itemize}
  \item Tema naj bo dovolj ozka, da jo je mogoče poglobljeno obdelati na 10--15 straneh.
  \item Vključi vsaj en izviren, samostojno rešen primer ali problem.
  \item Pokaži povezavo matematike z realnim svetom, drugo stroko ali zgodovino.
  \item Dokaze prilagodi ravni gimnazije: poudarek na razumevanju ideje, ne na formalni strogosti.
  \item Uporabi vizualizacije (grafi, geometrijske skice) povsod, kjer pomagajo razumevanju.
  \item Naloga naj se konča z lastnim zaključkom, ne le s ,,prepisanim`` povzetkom vira.
\end{itemize}

\subsubsection*{Okviren časovni potek}

\begin{enumerate}
  \item Izbira in potrditev teme pri mentorju.
  \item Zbiranje virov in literature.
  \item Izdelava grobega osnutka (struktura + jedro).
  \item Pregled pri mentorju, popravki.
  \item Izdelava končne različice.
  \item Priprava na ustno predstavitev/zagovor (če je predviden).
\end{enumerate}

\subsection{Predlogi naslovov po področjih}

\subsubsection*{Zgodovina matematike}
\begin{itemize}
  \item Zlati rez skozi zgodovino, naravo in umetnost
  \item Fibonaccijevo zaporedje in njegova uporaba v naravi
  \item Pitagorov izrek: različni dokazi skozi zgodovino
  \item Zgodovina in izračun števila $\pi$
  \item Nastanek neevklidskih geometrij
  \item Kriptografija skozi zgodovino -- od Cezarjeve šifre do RSA
\end{itemize}

\subsubsection*{Algebra}
\begin{itemize}
  \item Kompleksna števila in njihova uporaba v elektrotehniki
  \item Reševanje enačb 3.\ in 4.\ stopnje (Cardano, Ferrari)
  \item Matrike pri reševanju sistemov linearnih enačb
  \item Polinomi in Hornerjev algoritem
  \item Uvod v teorijo grup na preprostih primerih
\end{itemize}

\subsubsection*{Geometrija}
\begin{itemize}
  \item Platonska telesa in njihove lastnosti
  \item Fraktali -- geometrija narave
  \item Zlati rez v arhitekturi
  \item Trigonometrija in triangulacija v navigaciji/GPS
  \item Stožnice (elipsa, parabola, hiperbola) in njihova uporaba
  \item Neevklidske geometrije -- kaj se zgodi, če Evklidov 5.\ postulat ne velja
\end{itemize}

\subsubsection*{Matematična analiza}
\begin{itemize}
  \item Odvod pri reševanju optimizacijskih problemov (npr.\ najmanjša poraba materiala)
  \item Uporaba integralov pri izračunu površin in prostornin vrtenin
  \item Zenonovi paradoksi in pojem limite
  \item Diferencialne enačbe pri modeliranju rasti populacij
  \item Taylorjeva vrsta in aproksimacija funkcij
\end{itemize}

\subsubsection*{Verjetnost in statistika}
\begin{itemize}
  \item Verjetnost pri igrah na srečo (kocke, karte, loto)
  \item Statistična analiza podatkov (npr.\ rezultati dijakov, športni podatki)
  \item Bayesov izrek in njegova uporaba (npr.\ medicinsko testiranje)
  \item Slučajni sprehodi (random walk)
  \item Zakon velikih števil na simulacijah
\end{itemize}

\subsubsection*{Diskretna matematika in kombinatorika}
\begin{itemize}
  \item Teorija grafov in problem štirih barv
  \item Pascalov trikotnik in binomski izrek
  \item Osnove teorije iger
  \item Problem trgovskega potnika
\end{itemize}

\subsubsection*{Uporabna matematika}
\begin{itemize}
  \item Matematika v kriptografiji in varnosti podatkov
  \item Obrestno obrestovanje in matematika financ
  \item Matematika v računalniški grafiki
  \item Matematika glasbe -- frekvence, harmonija, intervali
  \item Matematični modeli širjenja epidemij (SIR model)
  \item Matematika pri gradnji mostov in zgradb
  \item GPS in matematika trilateracije
\end{itemize}

\subsubsection*{Teorija števil}
\begin{itemize}
  \item Praštevila in njihova vloga v kriptografiji
  \item Popolna in prijateljska števila
  \item Diofantske enačbe
  \item Modularna aritmetika in njena uporaba
\end{itemize}

\bigskip
\noindent\textit{Dijak se usmeri glede na njegov interes (npr.\ šport, glasba, računalništvo, ekonomija) -- tako je motivacija za samostojno delo večja, naloga pa bolj izvirna.}

\newpage
% =====================================================================
\section{Seminarska naloga}
\label{sec:seminarska}
% =====================================================================

Kratka seminarska naloga (do 3 strani, \LaTeX/Overleaf) -- strogo se drži naslova in uporablja samo formule, obravnavane v gimnaziji. Struktura: uvod (situacija, zgodovina/kontekst) -- jedro (matematično ozadje in konkreten izračun) -- zaključek (ugotovitve) -- literatura. Dijak lahko pri pisanju uporabi AI kot pomoč, vsebinsko pa mora ostati tesno vezan na izbrano temo in sklop snovi.

\medskip
\noindent Snov je navedena z oznako poglavja iz dokumenta \textit{Matematika na maturi -- cilji}, da lahko nalogo neposredno povežemo s standardi znanja.

\subsection{1.--2.\ letnik}

\subsubsection*{1. Naravna in cela števila}
\begin{itemize}
  \item \textbf{Usklajevanje izmen na dveh strojih/linijah} -- najmanjši skupni večkratnik  ciklov, kdaj se urnika znova poravnata.
  \item \textbf{Zobniki različnih velikosti} -- zakaj imajo zobniki pogosto praštevilsko število zob (enakomerna obraba), v/D pri vrtenju.
  \item \textbf{Kontrolna števka pri EAN/ISBN kodi} -- deljivost in modularni izračun kontrolne števke izdelka/knjige.
  \item \textbf{Shranjevanje podatkov v računalniku} -- pretvorba med desetiškim in dvojiškim zapisom (bit, bajt).
\end{itemize}

\subsubsection*{2. Algebrski izrazi, enačbe in neenačbe}
\begin{itemize}
  \item \textbf{Mešanje raztopin različnih koncentracij} -- npr.\ sirup, čistilo, kislina; linearna enačba za želeno mešanico.
  \item \textbf{Delitev računa/najema med prijatelji glede na porabo} -- postavitev in reševanje enačbe.
  \item \textbf{Prilagajanje recepta za drugo število oseb} -- sorazmerja in enačbe.
  \item \textbf{Načrtovanje žepnine/proračuna za izlet} -- neenačbe (koliko največ lahko porabim).
  \item \textbf{Vietovo pravilo pri hitrem preverjanju rešitev kvadratne enačbe} -- npr.\ pri nalogah o metu žoge (čas leta).
\end{itemize}

\subsubsection*{3. Osnove logike}
\begin{itemize}
  \item \textbf{Logična vezja v elektronskih napravah} -- AND, OR, NOT vezja in njihova uporaba.
  \item \textbf{Iskanje po spletu ali bazi podatkov z logičnimi operatorji} -- AND/OR/NOT pri iskalnih poizvedbah.
  \item \textbf{Uganke o resnicoljubcih in lažnivcih} -- logično sklepanje in izjavni račun.
  \item \textbf{Veljavnost pravnih/sodnih sklepov} -- osnovni silogizmi in pravilno sklepanje.
\end{itemize}

\subsubsection*{4. Množice}
\begin{itemize}
  \item \textbf{Vennovi diagrami pri šolski anketi} -- npr.\ dijaki, ki obiskujejo matematični IN športni krožek.
  \item \textbf{Filtri v spletni trgovini} -- presek in unija množic izdelkov.
  \item \textbf{Poizvedbe v Excelu/bazi podatkov} -- unija, presek, razlika podatkovnih naborov.
  \item \textbf{Alergije in sestavljanje jedilnika} -- dovoljena/prepovedana živila kot množice.
\end{itemize}

\subsubsection*{5. Racionalna števila}
\begin{itemize}
  \item \textbf{Popusti in akcije v trgovini} -- odstotki, promili, primerjava dejanskih prihrankov.
  \item \textbf{Obrestna mera pri potrošniškem kreditu} -- odstotni račun na konkretnem primeru ponudbe.
  \item \textbf{Statistika metov v košarki ali strelov na gol} -- odstotek uspešnosti.
  \item \textbf{Recepti in razredčevanje} -- barve, čistila; razmerja in deleži.
\end{itemize}

\subsubsection*{6. Realna števila}
\begin{itemize}
  \item \textbf{Zaokroževanje pri tehničnih meritvah} -- napake meritev v gradbeništvu ali strojništvu.
  \item \textbf{Dokaz iracionalnosti $\sqrt{2}$} -- Pitagorejci in diagonala kvadrata ter razlogi za propad pitagorejske šole.
  \item \textbf{Natančnost GPS koordinat} -- koliko decimalnih mest dejansko pomeni na terenu.
  \item \textbf{Merilna negotovost pri meritvi temperature ali dolžine} -- zapis z intervalom.
\end{itemize}

\subsubsection*{7. Pravokotni koordinatni sistem}
\begin{itemize}
  \item \textbf{Igra Bojne ladje ali šahovnica} -- koordinatni sistem v praksi.
  \item \textbf{Iskanje lokacij na zemljevidu/GPS} -- uporaba koordinat.
  \item \textbf{Zračna razdalja med dvema krajema} (npr.\ Krško--Ljubljana) -- izračun s formulo za razdaljo med točkama.
  \item \textbf{Piksli na zaslonu telefona/računalnika} -- koordinatni sistem digitalne slike.
\end{itemize}

\subsubsection*{8. Linearna funkcija}
\begin{itemize}
  \item \textbf{Cenik taksija ali parkirišča} -- začetna postavka + cena na km/uro kot linearna funkcija.
  \item \textbf{Poraba goriva glede na prevoženo razdaljo}.
  \item \textbf{Pretvorba temperature med Celzijem in Fahrenheitom} -- linearna zveza.
  \item \textbf{Najem opreme (kolesa, šotora, ozvočenja)} -- fiksni in spremenljivi stroški.
  \item \textbf{Načrtovanje proizvodnje z omejenimi viri} (npr.\ pekarna, čebelarstvo) -- linearno programiranje.
\end{itemize}

\subsubsection*{9. Evklidska geometrija v ravnini}
\begin{itemize}
  \item \textbf{Merjenje širine reke ali druge nedostopne razdalje} -- podobnost trikotnikov.
  \item \textbf{Načrtovanje vrta ali igrišča} -- obseg in ploščina lika.
  \item \textbf{Talesov izrek in obodni/središčni kot} -- pri načrtovanju krožnih objektov (rondo, kupola, kolo).
  \item \textbf{Zlati rez v oblikovanju} -- logotip, fotografija, embalaža.
\end{itemize}

\subsubsection*{10. Vektorji v ravnini in prostoru}
\begin{itemize}
  \item \textbf{Vlečenje sani ali vozička pod kotom} -- seštevanje sil kot vektorjev.
  \item \textbf{Navigacija čolna ali letala ob vetru/toku} -- vektorsko seštevanje hitrosti.
  \item \textbf{Premikanje lika v videoigri} -- vektor premika in hitrosti.
  \item \textbf{Smer in razdalja pri GPS navigaciji} -- zapis poti kot vektorja.
\end{itemize}

\subsubsection*{11. Potence z racionalnimi eksponenti}
\begin{itemize}
  \item \textbf{Škaliranje arhitekturnih maket} -- koreni pri preračunu površin/prostornin ob spremembi merila.
  \item \textbf{Čas prostega pada predmeta} -- koren v formuli za čas padca z dane višine.
  \item \textbf{Izračun stranice kvadrata ali polmera kroga iz znane ploščine} -- korenjenje.
\end{itemize}

\subsubsection*{12. Kompleksna števila}
\begin{itemize}
  \item \textbf{Izmenični tok v elektrotehniki} -- impedanca kot kompleksno število.
  \item \textbf{Fraktali (Mandelbrotova množica)} -- kompleksna ravnina in iteracija.
  \item \textbf{Nihanje/dušenje v fiziki} -- kvadratne enačbe brez realnih rešitev.
\end{itemize}

\subsubsection*{13. Funkcije in njene lastnosti}
\begin{itemize}
  \item \textbf{Cenik vozovnice ali parkirnine po conah} -- odsekoma podana funkcija.
  \item \textbf{Dnevni diagram porabe elektrike ali vode} -- branje lastnosti funkcije (naraščanje, ničle, ekstremi) iz grafa.
  \item \textbf{Pretvorbe med valutami ali merskimi enotami} -- funkcija in njen predpis.
\end{itemize}

\subsubsection*{14. Potenčna in korenska funkcija}
\begin{itemize}
  \item \textbf{Odvisnost prostornine/površine posode od dolžine roba} -- kvadratna/kubična rast.
  \item \textbf{Odpor zraka pri padalstvu} -- korenska odvisnost hitrosti padanja.
  \item \textbf{Keplerjev tretji zakon} -- odvisnost obhodnega časa planeta ($T$) od velike polosi tirnice ($a$): $T^2 \sim a^3$.
\end{itemize}

\subsubsection*{15. Kvadratna funkcija}
\begin{itemize}
  \item \textbf{Višina žoge v odvisnosti od časa pri metu} -- košarka, nogomet, odbojka.
  \item \textbf{Optimizacija ograjenega vrta pri dani dolžini ograje} -- največja možna ploščina.
  \item \textbf{Dobiček podjetja v odvisnosti od cene izdelka} -- kvadratni model dobička.
  \item \textbf{Oblika mostu ali loka} -- parabolični profil konstrukcije.
\end{itemize}

\subsection{3.--4.\ letnik}

\subsubsection*{16. Eksponentna in logaritemska funkcija}
\begin{itemize}
  \item \textbf{Rihterjeva lestvica merjenja potresov} -- logaritemska lestvica jakosti; primerjava potresov v Sloveniji (npr.\ Krško 1917, Posočje 1998/2004).
  \item \textbf{pH lestvica} -- logaritemska lestvica kislosti/bazičnosti; primeri iz vsakdana (voda, kisline v gospodinjstvu).
  \item \textbf{Decibelna lestvica jakosti zvoka} -- zakaj je zvok merjen logaritemsko; primerjava glasnosti prometa, koncerta, praga bolečine.
  \item \textbf{Obrestno obrestovanje} -- rast privarčevanega zneska/posojila; primerjava ponudb slovenskih bank.
  \item \textbf{Razpolovni čas radioaktivnih snovi} -- ogljikovo datiranje (C-14) pri arheoloških najdbah.
  \item \textbf{Eksponentna rast populacije bakterij/prebivalstva} -- modeliranje z eksponentno funkcijo.
  \item \textbf{Newtonov zakon ohlajanja} -- ohlajanje kave/telesa, eksponentni upad temperature.
  \item \textbf{Amortizacija vrednosti avtomobila} -- eksponentni padec vrednosti rabljenega vozila skozi leta.
  \item \textbf{Širjenje virusa v začetni fazi epidemije} -- eksponentni model rasti (npr.\ gripa, COVID-19 v prvih tednih).
\end{itemize}

\subsubsection*{17. Kotne funkcije}
\begin{itemize}
  \item \textbf{Merjenje višine objekta s kotnimi funkcijami} -- npr.\ dimnik TEB Brestanica, cerkveni zvonik, drevo; merjenje kota in razdalje ter izračun višine.
  \item \textbf{Dolžina sence skozi dan/leto} -- odvisnost od sončne elevacije, trigonometrični model.
  \item \textbf{Nihanje gugalnice ali nihala} -- modeliranje odmika s sinusno funkcijo.
  \item \textbf{Plimovanje} -- nihanje višine morja kot sinusna funkcija časa.
  \item \textbf{Izmenična električna napetost (AC)} -- sinusni signal v gospodinjskem omrežju (230~V, 50~Hz).
  \item \textbf{Koti med kazalci ure} -- izračun kota med urnim in minutnim kazalcem ob poljubni uri.
  \item \textbf{GPS in triangulacija} -- določanje lege s pomočjo kotov/razdalj do znanih točk.
\end{itemize}

\subsubsection*{18. Krivulje drugega reda (stožnice)}
\begin{itemize}
  \item \textbf{Parabolične satelitske antene} -- lastnost žarišča parabole in zbiranje signala.
  \item \textbf{Balistika -- let žoge ali izstrelka} -- parabolična trajektorija, doseg in največja višina.
  \item \textbf{Keplerjevi zakoni -- eliptične tirnice planetov} -- Zemljina tirnica kot elipsa.
  \item \textbf{Viseči most vs.\ verižnica in parabola} -- primerjava oblik nosilnih kablov (npr.\ na konkretnem mostu).
  \item \textbf{Šepetalnice (whispering galleries)} -- eliptična lastnost odboja zvoka v prostorih eliptičnega tlorisa.
  \item \textbf{Sončna peč / solarni koncentrator} -- uporaba parabole za zbiranje sončne energije.
\end{itemize}

\subsubsection*{19. Polinomi in racionalne funkcije}
\begin{itemize}
  \item \textbf{Optimalna prostornina škatle iz pole papirja} -- kubični polinom in iskanje maksimuma (lahko brez odvoda, z grafom/tabelo).
  \item \textbf{Modeliranje stroškov proizvodnje} -- polinomska funkcija skupnih stroškov v odvisnosti od količine.
  \item \textbf{Aproksimacija izmerjenih podatkov s polinomom} -- npr.\ dnevno nihanje temperature.
  \item \textbf{Zavorna pot vozila} -- odvisnost zavorne poti od hitrosti (kvadratni/polinomski model).
\end{itemize}

\subsubsection*{22. Zaporedja}
\begin{itemize}
  \item \textbf{Hanojski stolpi} -- število potez kot geometrijsko zaporedje/eksponentna rast.
  \item \textbf{Prepogibanje papirja} -- geometrijsko zaporedje debeline lista (in zakaj je 7-krat praktična meja).
  \item \textbf{Verižna pisma / piramidne sheme} -- geometrijska rast števila udeležencev.
  \item \textbf{Anuitetno varčevanje ali odplačevanje kredita} -- aritmetično/geometrijsko zaporedje obrokov.
  \item \textbf{Fibonaccijevo zaporedje v naravi} -- storžki, sončnice, razporeditev listov.
\end{itemize}

\subsubsection*{23. Diferencialni račun}
\begin{itemize}
  \item \textbf{Optimizacija embalaže} -- pločevinka ali škatla z najmanjšo porabo materiala pri dani prostornini.
  \item \textbf{Hitrost in pospešek vozila} -- odvod poti po času iz podatkov o vožnji (npr.\ GPS sledenje).
  \item \textbf{Mejni stroški in mejni dobiček} -- odvod funkcije stroškov/prihodka v ekonomiji, optimalna cena izdelka.
  \item \textbf{Reševalec na plaži (Snellov problem)} -- najhitrejša pot do tonočega ob spremembi hitrosti (voda/pesek).
  \item \textbf{Oblika mostu ali loka za minimalno porabo materiala} -- optimizacijski problem z odvodom.
\end{itemize}

\subsubsection*{24. Integralni račun}
\begin{itemize}
  \item \textbf{Prevožena pot iz grafa hitrosti} -- integral hitrosti po času (npr.\ iz podatkov vozila ali kolesarjenja).
  \item \textbf{Površina nepravilne parcele ali jezera} -- numerična ocena s pomočjo integrala/vsote pravokotnikov.
  \item \textbf{Prostornina vrtenine} -- kozarec, sod ali vaza kot telo, dobljeno z vrtenjem krivulje.
  \item \textbf{Delo raztezanja vzmeti (Hookov zakon)} -- izračun opravljenega dela z integralom sile.
  \item \textbf{Poraba goriva iz diagrama hitrosti/pretoka} -- skupna poraba kot ploščina pod grafom.
\end{itemize}

\subsubsection*{25. Kombinatorika}
\begin{itemize}
  \item \textbf{Varnost gesel in PIN kod} -- število možnih kombinacij glede na dolžino in nabor znakov.
  \item \textbf{Žreb športne lige ali turnirja} -- število možnih razporedov/parov tekem.
  \item \textbf{Registrske tablice} -- koliko različnih kombinacij dopušča slovenski sistem tablic.
  \item \textbf{Loto 7 od 39} -- verjetnost zadetka glede na število možnih kombinacij.
\end{itemize}

\subsubsection*{26. Verjetnostni račun}
\begin{itemize}
  \item \textbf{Verjetnost pri igrah na srečo} -- kocka, karte, ruleta.
  \item \textbf{Medicinsko testiranje in Bayesov izrek} -- verjetnost dejanske bolezni ob pozitivnem testu (lažno pozitivni rezultati).
  \item \textbf{Rojstnodnevni paradoks} -- verjetnost, da imata dva dijaka v razredu isti rojstni dan.
  \item \textbf{Zavarovalna premija glede na tveganje} -- osnovni izračun pričakovane vrednosti izplačila.
  \item \textbf{Vremenska napoved -- verjetnost padavin} -- kako razumeti in interpretirati odstotek verjetnosti.
\end{itemize}

\subsubsection*{27. Statistika}
\begin{itemize}
  \item \textbf{Analiza uspeha dijakov pri maturi} -- statistična obdelava (povprečje, mediana, standardni odklon) realnih ali javno dostopnih podatkov.
  \item \textbf{Anketa med dijaki šole} -- zbiranje in statistična obdelava lastnih podatkov (npr.\ čas za domačo nalogo, raba telefona).
  \item \textbf{Statistika prometnih nesreč v Posavju} -- analiza javno dostopnih podatkov (npr.\ AVP, SURS).
  \item \textbf{Športna statistika} -- analiza rezultatov kluba ali reprezentance (povprečja, odstopanja).
  \item \textbf{Poraba elektrike/vode v gospodinjstvu čez leto} -- statistična analiza in iskanje trenda.
\end{itemize}

\bigskip
\noindent\textit{Namig: pri izbiri naj dijak preveri, da za izbrano situacijo obstajajo dostopni realni podatki (uradne statistike, meritve, javno dostopni viri), saj to nalogo naredi bolj konkretno in verodostojno.}

\end{document}
